% Zone „Das kennst du schon“ – aus bank/symmetrie-abbildungen/zone.jsonl (zusammenbau v0.1)
\einheitenkopf[zone][Kennst du schon]{Das kennst du schon}
\setcounter{aufgabe}{0}
%% TODO Grundvorstellungs-Aufgabe als erste Hauptnummer der Zone (2.8) fehlt in der Bank

%% TODO Titel: Ich-kann-Satz fehlt in der Bank (Platzhalter: Kettenname) – Nr. 1
%% TODO Anweisung: ein Satz über der ersten Teilaufgabe fehlt in der Bank – Nr. 1
\begin{aufgabe}{Kästchen im Raster abzählen und Strecken auf dem Karo abtragen}
%% TODO Darstellung für schwach fehlt in der Bank (symmetrie-abbildungen-zone-f1-v1; Abschnitt „Für schwache Schüler“)
\swz[3]{Ein Kästchen auf dem Karopapier ist 5\,mm breit. Eine Strecke ist 4 Kästchen lang. Wie viele Millimeter?}{}
%% TODO Darstellung für schwach fehlt in der Bank (symmetrie-abbildungen-zone-f1-v3; Abschnitt „Für schwache Schüler“)
\swz[3]{Eine Strecke auf dem Karo ist 4{,}5\,cm lang, ein Kästchen 5\,mm breit. Wie viele Kästchen?}{}
\end{aufgabe}

%% TODO Titel: Ich-kann-Satz fehlt in der Bank (Platzhalter: Kettenname) – Nr. 2
%% TODO Anweisung: ein Satz über der ersten Teilaufgabe fehlt in der Bank – Nr. 2
\begin{aufgabe}{Senkrechte und parallele Geraden mit dem Geodreieck zeichnen und erkennen, Abstand eines Punktes von einer Geraden messen}
\begin{teile}
\teil Stehen die beiden Geraden im Bild senkrecht zueinander? Prüfe mit dem Geodreieck. \\ \kreuz{ja} \\ \kreuz{nein}
\end{teile}
%% TODO Darstellung für schwach fehlt in der Bank (symmetrie-abbildungen-zone-f2-v3; Abschnitt „Für schwache Schüler“)
\swz[3]{Von einem Punkt P führen drei Strecken zur Geraden g: 4{,}2\,cm, 3{,}6\,cm und 5{,}1\,cm lang. Eine davon steht senkrecht auf g. Wie groß ist der Abstand von P zu g?}{}
\end{aufgabe}

%% TODO Titel: Ich-kann-Satz fehlt in der Bank (Platzhalter: Kettenname) – Nr. 3
%% TODO Anweisung: ein Satz über der ersten Teilaufgabe fehlt in der Bank – Nr. 3
\begin{aufgabe}{Negative Zahlen an der Zahlengeraden ablesen und eintragen}
\swz[3]{Welche Zahl gehört zum Punkt A?}{\zahlenstrahl[xmin=-5,xmax=5,xstep=1,karo=0.8]{-3/A, 2/B}}
\swz[3]{Welche Zahl gehört zum Punkt D?}{\zahlenstrahl[xmin=-5,xmax=5,xstep=1,karo=0.8]{-2.5/D, 1/}}
\end{aufgabe}

%% TODO Titel: Ich-kann-Satz fehlt in der Bank (Platzhalter: Kettenname) – Nr. 4
%% TODO Anweisung: ein Satz über der ersten Teilaufgabe fehlt in der Bank – Nr. 4
\begin{aufgabe}{Vierecksarten und Dreiecksarten an ihren Eigenschaften erkennen und benennen (Quadrat, Rechteck, Raute, Parallelogramm, Drachenviereck, gleichschenkliges Trapez; gleichseitiges und gleichschenkliges Dreieck)}
\swz[3]{Wie heißt dieses Viereck?}{\rechteck{4}{2}}
%% TODO Darstellung für schwach fehlt in der Bank (symmetrie-abbildungen-zone-f4-v3; Abschnitt „Für schwache Schüler“)
\swz[3]{Ein Viereck hat vier gleich lange Seiten, aber keinen rechten Winkel. Wie heißt es?}{}
\end{aufgabe}

%% TODO Titel: Ich-kann-Satz fehlt in der Bank (Platzhalter: Kettenname) – Nr. 5
%% TODO Anweisung: ein Satz über der ersten Teilaufgabe fehlt in der Bank – Nr. 5
\begin{aufgabe}{Diagonalen eines Vierecks einzeichnen und benennen}
\swa{Zeichne im Viereck ABCD die Diagonale ein, die bei A beginnt. Wie heißt sie?}{\viereck{(0,0)}{(4,0)}{(5,3)}{(1,3.5)}}
\swa{Zeichne im Viereck ABCD beide Diagonalen ein und benenne sie.}{\viereck{(0,0)}{(5,1)}{(4,4)}{(0.5,3)}}
\end{aufgabe}

%% TODO Titel: Ich-kann-Satz fehlt in der Bank (Platzhalter: Kettenname) – Nr. 6
%% TODO Anweisung: ein Satz über der ersten Teilaufgabe fehlt in der Bank – Nr. 6
\begin{aufgabe}{Winkel messen und zeichnen, rechten Winkel erkennen und markieren}
\swz[3]{Miss den Winkel $\alpha$.}{\winkel{60}{\alpha}}
\swa{Zeichne an den Strahl bei S einen Winkel von $125^\circ$.}{\winkelstrahl}
\end{aufgabe}

%% TODO Titel: Ich-kann-Satz fehlt in der Bank (Platzhalter: Kettenname) – Nr. 7
%% TODO Anweisung: ein Satz über der ersten Teilaufgabe fehlt in der Bank – Nr. 7
\begin{aufgabe}{Strecken messen und gleich lange Strecken abtragen}
\swz[3]{Miss die längste Seite c des Dreiecks.}{\dreieckrw{4}{3}{a}{b}{c}}
\swz[3]{Miss die Seite c auf den Millimeter genau.}{\dreieckrw{3.6}{2.7}{a}{b}{c}}
\end{aufgabe}

%% TODO Titel: Ich-kann-Satz fehlt in der Bank (Platzhalter: Kettenname) – Nr. 8
%% TODO Anweisung: ein Satz über der ersten Teilaufgabe fehlt in der Bank – Nr. 8
\begin{aufgabe}{Kästchen im Raster abzählen und Strecken auf dem Karo abtragen – Fehler finden}
\swz[3]{Finde den Fehler. Jonas soll abzählen, wie viele Kästchen eine Strecke von der 3. bis zur 8. Karolinie lang ist. Er zählt die Linien: 3, 4, 5, 6, 7, 8 – also 6 Kästchen.}{}
\end{aufgabe}

%% TODO Titel: Ich-kann-Satz fehlt in der Bank (Platzhalter: Kettenname) – Nr. 9
%% TODO Anweisung: ein Satz über der ersten Teilaufgabe fehlt in der Bank – Nr. 9
\begin{aufgabe}{Kästchen im Raster abzählen und Strecken auf dem Karo abtragen – selbst rechnen}
%% TODO Darstellung für schwach fehlt in der Bank (symmetrie-abbildungen-zone-f1-v6; Abschnitt „Für schwache Schüler“)
\swz[3]{Eine Strecke reicht von der 4. bis zur 11. Karolinie. Wie viele Kästchen lang?}{}
\end{aufgabe}
